% Implementation of the pricing test: the statistic, the boost, the scoring
% path and the status of the preconditions (Table tab:preconditions), moved
% from or cited by Section~\ref{sec:pricing-test}. No heading: the integrator
% inputs this file inside Appendix~B (Experimental details), e.g. at the end
% of the subsection on metric implementations. \label{app:scoring} takes the
% number of the enclosing (sub)section and is cited from pricing.tex.
\label{app:scoring}%
The statistic of the pricing test, Eq.~\eqref{eq:statistic}, is computed on
the luminance
$y = 0.299\,x_{\mathrm{R}} + 0.587\,x_{\mathrm{G}} + 0.114\,x_{\mathrm{B}}$ of
an RGB image $x$ with values in $[0, 255]$ (ITU-R BT.601 weights, unrounded).
Its 2-D DFT $\hat y = F y$ on the $h \times w$ pixel grid is indexed by $u$ after the
zero frequency is shifted to $c = (\lfloor h/2 \rfloor, \lfloor w/2 \rfloor)$;
the integer radius is $r(u) = \lfloor \lVert u - c \rVert_2 \rfloor$, and the
bins are $U_k = \{u : r(u) = k\}$, $k = 0, \dots, K-1$. Averaging decibels
rather than power makes $s$ a geometric-mean spectrum, and the normalisation of
$F$ cancels in every difference used. Because $s$ sees only the luminance, the
operators impose its change on every colour channel, and a change in chroma is
neither measured nor imposed separately. The boost is the gain
$g = 10^{d \varrho / 20}$ of dose $d$ in decibels, with the raised-cosine ramp
$\varrho = \frac12 - \frac12 \cos\big(\pi \min\{1, \max\{0, (\nu - 0.667)/0.5\}\}\big)$
in the radius $\nu = \lVert u - c \rVert_2 / (\min(h, w)/2)$, measured in units
of the axial Nyquist frequency. White noise $\xi$ is drawn independently for
every pixel and channel and added without rounding.

Every term of Eqs.~\eqref{eq:value}--\eqref{eq:accounting} is scored in
floating point. The filtered images, clipped to $[0, 255]$ and not rounded,
enter the `clean' resizer and feature extractor of \texttt{clean-fid}
\citep{parmar2022aliased} directly: bicubic resizing of each channel to
$299 \times 299$ in floating point, 2048-dimensional Inception-V3 pool
features \citep{szegedy2016rethinking}, and the Fr\'echet distance to the
reference features \citep{heusel2017gans}. FD$_{\text{DINOv2}}$
(Section~\ref{sec:worth-dino}) replaces the extractor by DINOv2 ViT-L/14
\citep{oquab2024dinov2}: images are resized bicubically to $224 \times 224$ in
floating point, normalised with the ImageNet mean and standard deviation, and
represented by the model's normalised class token, as in
\citet{stein2023exposing} except that the resize acts on floating-point rather
than 8-bit images, for original and filtered samples alike. For unmodified samples the path
reproduces \texttt{clean-fid}'s file-based computation: the baseline FIDs
agree with those of the training ledger to within \priceIdentityMaxDiff. The
identity operator used for this check bypasses the filter and the clipping,
so the filter path at zero change is exercised only by arms whose change is
close to zero, such as the static weighting of the null control. Runs that
reproduce an earlier 8-bit measurement round the filtered images, with or
without a dither of $\pm 0.5$~LSB, to isolate the share of the value due to
re-quantisation (hazard H1). The test costs one computation of $s$ per set
and one feature-extraction pass per operator and seed.

Two effects of the 8-bit input remain in the floating-point values. The
operators act on $B$'s stored 8-bit samples, so a gain $g$ also rescales $B$'s
own rounding-noise field; in each bin this is equivalent, in power, to
adding independent noise of $g^2 - 1$ times the rounding variance, a field of
the kind the white-noise control prices. Its share of $V$ is not separated, and
it grows with $g^2 - 1$, most for the larger boost doses and the match.
Clipping, in turn, makes the change measured back after an operator differ
slightly from the imposed $\delta_i$, most for arms whose change is large;
Table~\ref{tab:pricing-arms} reports the change measured back.

Table~\ref{tab:preconditions} lists where each precondition of
Section~\ref{sec:pricing-pre} is checked and whether it holds in the three
settings.

\begin{table}[tbp]
\centering
\small
\caption{Status of the preconditions of the pricing test. Section: where the
  CIFAR-10 evidence is; that for CelebA-64 is in Section~\ref{sec:celeba}, that
  for DreamBooth in Section~\ref{sec:dreambooth}. P3: on CIFAR-10 every
  spectral operator is a coherent radial gain and only the transplant imposes
  $\delta_A$; on CelebA-64 the coherent gains and white noise move FID in
  opposite directions. n/a: needs the AWWL
  arm, which was not trained on CelebA-64; not run: the check was not
  performed.}
\label{tab:preconditions}
\begin{tabular}{@{}llllll@{}}
\toprule
 & Requirement & Section & CIFAR-10 & CelebA-64 & DreamBooth \\
\midrule
P0 & no re-quantisation & \ref{sec:worth-curve} & met & met
  & not met (8-bit) \\
P1 & $s$ moves as scored & \ref{sec:mechanism} & DDIM-100 only & n/a
  & met \\
P2 & $\delta$ measured back & \ref{sec:worth-headline} & met (transplant)
  & n/a & not run \\
P3 & several operators & \ref{sec:worth-checks}
  & coherent gains only & signs differ & boost only \\
P4 & metric registers $s$ & \ref{sec:worth-curve}
  & met & met & not met \\
P5 & no interaction & \ref{sec:worth-checks} & consistent & n/a
  & not run \\
\bottomrule
\end{tabular}
\end{table}
